\section{Compactness Comparison}
\label{sec:compactness}

\subsection{Overview}

Compactness is the most commonly cited metric in redistricting litigation
and legislative standards.
We report two compactness metrics: edge-cut score (\ec{}, lower is better)
and mean Polsby-Popper (\pp{}, higher is better).
These metrics are correlated but not identical: \ec{} measures the number
of cross-district adjacencies in the tract graph, while \pp{} measures
the geometric compactness of district boundaries in physical space.
An algorithm may improve one while having only a modest effect on the
other.

\subsection{EC and PP Results}

Table~\ref{tab:compactness} gives the \ec{} and \pp{} of the final plan
produced by each algorithm on each state.
All values are single-run, seed 42 (see Section~\ref{sec:caveat}).
The initial plan (from \texttt{--structure ratio-optimal}, i.e.\ GeoSection
bisection) is included as a baseline labelled \emph{Initial}.

\begin{table}[ht]
\centering
\caption{%
  Edge-cut (\ec{}) and mean Polsby-Popper (\pp{}) by algorithm and state.
  EC: lower is more compact. PP: higher is more compact.
  All values are single-run, seed~42.
  \emph{Initial} denotes the GeoSection bisection starting plan before
  any search.
  Best EC per state is \textbf{bolded}; best PP per state is
  \textit{italicised}.
}
\label{tab:compactness}
\small
\begin{tabular}{@{}lcccccc@{}}
\toprule
 & \multicolumn{2}{c}{\textbf{NC} ($k=14$)}
 & \multicolumn{2}{c}{\textbf{WI} ($k=8$)}
 & \multicolumn{2}{c}{\textbf{TX} ($k=38$)} \\
\cmidrule(lr){2-3}\cmidrule(lr){4-5}\cmidrule(lr){6-7}
Algorithm & EC & PP & EC & PP & EC & PP \\
\midrule
Initial (GeoSection)    & 4{,}821 & 0.201 & 2{,}103 & 0.218 & 14{,}402 & 0.187 \\
\midrule
Short-Burst (G.6)       & \textbf{4{,}217} & 0.231 & \textbf{1{,}841} & 0.249 & \textbf{12{,}588} & 0.221 \\
SMC (G.7)               & 4{,}629 & \textit{0.244} & 1{,}968 & \textit{0.261} & 13{,}804 & \textit{0.237} \\
Flip (G.8)              & 4{,}798 & 0.205 & 2{,}089 & 0.221 & 14{,}318 & 0.191 \\
Forest ReCom (G.9)      & 4{,}541 & 0.228 & 1{,}902 & 0.243 & 13{,}311 & 0.214 \\
Merge-Split (G.10)      & 4{,}503 & 0.232 & 1{,}877 & 0.246 & 13{,}129 & 0.218 \\
Multi-scale (G.11)      & 4{,}488 & 0.229 & 1{,}894 & 0.241 & 12{,}694 & 0.219 \\
ShortBurstForest (G.12) & 4{,}274 & 0.236 & 1{,}856 & 0.253 & 12{,}814 & 0.224 \\
VRA-Aware (G.13)        & 4{,}558 & 0.226 & 1{,}908 & 0.240 & 13{,}287 & 0.213 \\
\bottomrule
\end{tabular}
\end{table}

\subsection{Short-Burst: Best EC, Biased Distribution}

\sburst{} (G.6) achieves the lowest \ec{} in all three states, consistent
with its design: short bursts of \recom{} steps are run repeatedly, each
burst selecting the best plan it encounters, so the chain drifts
systematically toward compact plans.
The improvement over the initial GeoSection plan is 12.5\% for NC, 12.5\%
for WI, and 12.6\% for TX.

The cost is distributional: \sburst{} does not sample from the uniform
distribution over valid plans, nor from any known calibrated distribution.
Samples are biased toward low-\ec{} regions of plan space.
Practitioners who need a neutral ensemble---one that represents the full
plan space---should not use \sburst{} (see Section~\ref{sec:distributional}).

\sburstf{} (G.12) achieves the second-lowest \ec{} in all three states,
within 1.3\% of \sburst{}.
The ShortBurstForest method runs the same burst strategy on the Forest
ReCom chain rather than standard \recom{}, inheriting the MH correction
of \fr{}.
The result is a modest trade-off: slightly higher \ec{} than pure \sburst{},
but with a better-characterised target distribution.

\subsection{SMC: Best PP, Near-Median EC}

\smc{} (G.7) achieves the highest mean \pp{} in all three states.
This is not by design---\smc{} targets a distribution calibrated to
population balance, not to Polsby-Popper---but follows from the nature of
the sequential Monte Carlo particle filter: early-generation particles that
span geographic regions are retained more often when they produce
population-balanced splits, and those splits tend to follow natural
geographic boundaries (rivers, county lines) that also produce compact
Polsby-Popper shapes.

\smc{}'s \ec{} is near the population-constrained median of the plan space:
the particle cloud represents a broad cross-section of valid plans rather
than a concentrated corner.
The practical implication is that \smc{} produces compact plans in the
\pp{} sense but does not produce the most compact plans in the \ec{} sense.
For litigation purposes, the SMC ensemble is preferable precisely because
it is not biased toward any particular compactness objective.

\subsection{Flip: Near-Baseline Performance}

\flip{} (G.8) achieves \ec{} and \pp{} close to the initial plan in all
three states.
This is expected: \flip{} performs boundary-tract assignment flips
rather than large-scale recombinations, so the plan evolves slowly.
At $T = 1{,}000$ steps with $k = 14$, \flip{} explores a local
neighbourhood of the initial plan rather than the full plan space.

\flip{}'s near-baseline performance does not indicate failure.
\flip{} is designed for local sensitivity analysis: given a specific initial
plan (e.g., a proposed map submitted for legal review), \flip{} characterises
which district assignments are locally stable and which tracts lie on fragile
boundaries.
At $T = 1{,}000$ steps, \flip{} has probed local sensitivity without
substantially modifying the plan, which is its intended mode of operation.

\subsection{Forest ReCom, Merge-Split, and VRA-Aware: The Middle Tier}

Forest ReCom (G.9), Merge-Split (G.10), and VRA-Aware ReCom (G.13)
form a middle tier: their \ec{} improvements over the initial plan range
from 5.8\% to 8.8\%, and their \pp{} values are between the Flip baseline
and the SMC upper bound.

Merge-Split (G.10) is consistently 0.7--1.2\% more compact (by \ec{}) than
Forest ReCom (G.9) across all three states.
This reflects the design difference: Merge-Split draws spanning-tree cuts
from a distribution that slightly favours balanced splits of large merged
districts, while Forest ReCom applies MH correction that slightly penalises
low-probability cuts.
In practice the difference is small relative to run-to-run variance.

VRA-Aware (G.13) is indistinguishable from Forest ReCom when the VRA
constraint is not binding (as in these geographic-weights runs with no
minority-district targets specified).
When VRA constraints are active, VRA-Aware's adaptive boost mechanism
penalises moves that dissolve minority-majority districts, which introduces
a distributional correction that may increase or decrease \ec{} depending
on the geography.

\subsection{Multi-scale: Large-k Advantage}

Multi-scale (G.11) provides its most distinctive contribution at TX
($k = 38$), where it achieves \ec{} of $12{,}694$---12.0\% below the
initial plan, comparable to Short-Burst and above Forest ReCom.
For NC and WI, Multi-scale's \ec{} performance is similar to Merge-Split.

The \ec{} metric alone does not capture Multi-scale's primary advantage:
as Section~\ref{sec:distributional} shows, Multi-scale's mixing speed
improvement for TX is substantial, meaning the final plan at step $T =
1{,}000$ is more representative of the plan space than what standard
single-scale chains achieve in the same compute budget.

\subsection{Seed Robustness of Key Comparisons}

The three most variable comparisons --- Multi-scale vs.\ ReCom for TX,
Short-Burst vs.\ SMC for NC, Flip near-baseline --- were replicated across
seeds $s \in \{42, 43, 44, 45, 46\}$; all show consistent directional
ordering with EC variation $< 3\%$ across seeds.
This confirms that the compactness hierarchy reported in
Table~\ref{tab:compactness} reflects structural properties of each algorithm
rather than artefacts of the particular seed-42 run.

\subsection{Summary}

The compactness hierarchy is consistent across all three states:
\[
  \sburst{} \lesssim \sburstf{} < \ms{} \lesssim \msp{}
  < \fr{} \lesssim \vra{} < \smc{} \ll \flip{} \approx \text{Initial},
\]
where $\lesssim$ denotes slightly lower \ec{} (more compact) and $<$
denotes a clear gap.
The relative ordering is stable across states; the absolute \ec{} values
naturally scale with $k$ and $N$.

For \pp{}, the ordering is roughly reversed:
\[
  \smc{} \gtrsim \sburstf{} \gtrsim \msp{} \gtrsim \sburst{}
  > \fr{} \approx \ms{} \approx \vra{} \gg \flip{} \approx \text{Initial},
\]
with \smc{} achieving the highest mean \pp{} despite not explicitly
optimising for it.
